\section{Materials and Methods}
\label{sec:methods}

\subsection{Study design and data source}

We conducted a retrospective multicenter study using electronic health record
data from six Israeli general hospitals, drawn from Kineret,\citep{rappoport2025kineret} the data warehouse of Israel's largest hospital
network, represented in the Observational Medical Outcomes Partnership (OMOP)
common data model. The assessed recommendations are those of the national
position paper on the treatment of diabetes in hospitalization in a general
hospital,\citep{hochberg2023inpatient} issued in September 2023. The
pre-intervention period ends immediately before publication (August 2022--July
2023) and the post-intervention period begins a year after it (August
2024--December 2025), the intervening months being excluded to allow for
staggered adoption.

Kineret was not mapped to an international coding system at the time of this
work, so a dataset-specific concept index---a many-to-one mapping from local
codes onto the concepts named in the knowledge base---was curated first. That
index is what makes the remainder of the pipeline dataset-agnostic: once it
exists, the same knowledge base and engine run unchanged on any source within
the same clinical scope. The curated indices are provided as supplementary
material, and the curation procedure is described in Supplementary
Appendix~S1.

\subsection{Cohort definition}
\label{sec:cohort-definition}

An admission was eligible if it was an inpatient visit of 48 hours to 14 days
in a patient aged at least 18 years, with evidence of a genuine overnight stay,
no death within 48 hours of admission, no admission to an excluded department,
and at least one criterion for diabetes or hyperglycemia. Excluding pediatric
and obstetric departments also removes pregnancy-related admissions. Diabetes or
hyperglycemia was established by at least one of three non-mutually-exclusive
criteria: a documented diabetes
diagnosis from a manually curated list of concept identifiers; recurrent
hyperglycemia, two glucose measurements $\geq$180 mg/dL separated by 2 to 24
hours; or extreme hyperglycemia, any glucose measurement $\geq$250 mg/dL.
Hypoglycemia was deliberately not an entry criterion: the cohort is defined by
hyperglycemic burden, and hypoglycemia is treated throughout as an outcome
rather than as a route into the study population. Full eligibility rules are
given in Supplementary Appendix~S1.

\subsection{Abstraction input}

The \emph{temporal input} is a long-format table in which each row records one
time-stamped clinical fact, with structural tokens inserted to make the
admission timeline explicit---notably recurring \textsc{meal} tokens, assumed
to fall at fixed daily times because meal times are not recorded in the source
data. The \emph{context input} holds one row per admission of baseline
characteristics from the first 48 hours, from which a modified Charlson
comorbidity index (CCI) was computed (Supplementary Appendix~S1).

\subsection{Guideline adherence assessment}

Adherence was assessed by knowledge-based temporal abstraction,\citep{shahar1997framework} which converts time-stamped clinical data into
named interval concepts and then detects, over those intervals, the temporal
pattern that correct execution of each recommendation would have left in the
record.\citep{ben2024supporting} Abstraction and assessment are performed in a
single pass by our reimplementation of the Temporal Mediator.\citep{mediator-github} All domain knowledge is declared externally as XML
units validated against a shared schema, so that the clinical semantics and the
guideline requirements are stated once, outside the analysis code.

Each recommendation is formalized as a quality-assessment temporal pattern: an
\emph{anchor} identifying the clinical situation in which it applies, an
\emph{action} that satisfies it, a constraint type, and a membership function
grading partial compliance. Applicability is enforced by context---an anchor
falling outside the context in which the recommendation holds generates no
obligation and is not scored---so that an insulin dose after a meal is required
only where glucose is elevated. Recurring recommendations such as routine
glucose monitoring are instead evaluated by counting qualifying events within
fixed-length windows.

Compliance is graded rather than binary. Each constraint carries a trapezoidal
membership function mapping observed timing, frequency, or value to $[0,1]$: if
correction insulin is expected within 4 hours of a hyperglycemic reading and
considered missed after 12, insulin given at 8 hours scores 0.5 rather than
counting as an outright failure. Three constraint types are
used---\emph{time}, \emph{value}, and \emph{cyclic}---corresponding to the
compliance aspects reported below. An unanswered anchor
scores 0 where the recommendation requires an action, and 1 where it requires
one to be withheld, since there the absence of a response is itself compliant;
a recommendation that never became applicable yields no output rather than a
zero, so inapplicable recommendations do not depress the score.

Fourteen recommendations were formalized by the first author under the guidance
of two co-authors specializing in temporal reasoning and in diabetes care, who
reviewed them for clinical fidelity against the source position paper. Scoring
is retrospective and offline: the procedure issues no recommendation and
returns nothing to the treating team. Thresholds and the full pattern
definitions are given in Supplementary Appendices~S1 and~S3.

\subsection{Outcome definitions}

Ten complications were evaluated: 30-day mortality, hyperglycemia, severe
hyperglycemia, hypoglycemia, severe hypoglycemia, hyperosmolality,
ketoacidosis, infection, inpatient acute kidney injury, and inpatient
myocardial infarction. All ten are emitted as abstractions by the same engine,
so that outcome ascertainment and adherence assessment draw on a single
explicit knowledge base. Thirty-day mortality was taken from linked death
records, defined as death between 14 days before and 30 days after the recorded
discharge datetime, the backward tolerance capturing in-hospital deaths whose
death record precedes the discharge timestamp. Operational definitions are
given in Supplementary Appendix~S2.

\subsection{Statistical analysis}

The admission is the unit of analysis. Compliance scores were collapsed
successively over compliance aspects, calendar days, and finally patterns with
equal weight, so that neither frequently-evaluated patterns nor
frequently-evaluated days dominate the resulting per-admission score
(Supplementary Appendix~S1).

Adherence was compared between periods by Welch's unequal-variance $t$-test,
each admission attributed by its admission date. Confidence intervals are
unadjusted Welch $t$-intervals, and $P$ values were adjusted across
component-level comparisons by the Benjamini--Hochberg false discovery rate
(FDR) procedure.\citep{benjamini1995controlling} Monthly means carry 95\% CIs
from the within-month standard error, and the temporal trend is an ordinary
least-squares fit to those means.

Whole-admission adherence was compared between admissions ending in death
within 30 days and those not, and the graded relationship tested by logistic
regression on the continuous score, reported as the odds ratio (OR) per
0.1-point increase, then adjusted for age, sex, modified CCI, and hospital with
standard errors clustered on patient identity. Both are descriptive, since
whole-admission adherence does not strictly precede the outcome.

To identify which individual recommendations were associated with each
complication, one model was fitted per complication at cutoffs of 24, 48, and
72 hours. Each row is an admission--complication pair, labelled positive if the
complication first occurred at or after the cutoff and negative if it never
occurred; admissions in which it occurred \emph{before} the cutoff were
excluded, since a pre-existing condition cannot be attributed to subsequent
in-stay care, and adherence and covariates were computed only up to the cutoff.
Models were fitted by bias-reduced (Firth penalized-likelihood) logistic
regression,\citep{firth1993bias} unpenalized maximum likelihood having failed
to converge under separation for most complications. Correlated predictors were
pruned and the remainder standardized, and a predictor was considered robust
only if it remained significant with consistent direction across all three
windows. To avoid mechanical circularity, routine glucose monitoring was
excluded from the glycemic models, since glucose measurement is itself the
mechanism by which those events are detected.

Because these models estimate conditional predictive associations rather than
treatment effects, overall adherence and the recurrent components were
re-evaluated by inverse probability of treatment weighting, with adherence
dichotomized at the median of each outcome's eligible population. The binary
exposure is intrinsic to the design rather than a simplification, since the
propensity score \citep{rosenbaum1983central} and the standardized mean
differences (SMDs) used to verify balance \citep{austin2009balance} are
two-group constructs undefined for a continuous exposure; retaining continuous
adherence in the regressions above therefore makes the two designs
complementary. Propensity scores were estimated by bias-reduced logistic
regression, converted to stabilized weights and trimmed at the 1st and 99th
percentiles. Balance was assessed by the maximum post-weighting SMD, with
$<$0.10 required for an estimate to be interpretable; pairs failing this are
reported as non-estimable. Confidence intervals are unadjusted throughout, so an
interval excluding the null may accompany a nonsignificant adjusted $P$ value.
Analyses used Python 3.13 (Supplementary Appendix~S1).
